\documentclass[11pt]{article}
\usepackage{amsmath,amssymb,amsthm,geometry}
\geometry{margin=1in}
\newtheorem{thm}{Theorem}[section]
\newtheorem{lem}{Lemma}[section]
\newtheorem{prop}{Proposition}[section]
\newtheorem{rem}{Remark}[section]
\newcommand{\R}{\mathbb{R}}
\newcommand{\Om}{\sqrt{\kappa^{2}+\tau^{2}}}

\title{Constant Curvature and Constant Torsion Implies a Circular Helix\\
\large A self-contained differential-geometric proof}
\author{}
\date{}

\begin{document}
\maketitle

\begin{abstract}
We give a complete, fully explicit proof that a regular $C^{3}$ curve in
$\R^{3}$ with constant non-zero curvature $\kappa$ and constant torsion
$\tau$ is, up to a Euclidean motion, a circular helix. The degenerate cases
(circle, line) are classified as well. All derivatives, the curvature, and the
torsion are computed in closed form for the standard parametrisation
$\mathbf r(t)=(a\cos\omega t,a\sin\omega t,ct)$, and the recovery of the curve
from its Frenet--Serret data is carried out with correct bookkeeping of the
axial component. No physical interpretation is imposed; this is pure
differential geometry.
\end{abstract}

\section{Regularity and derivatives}
\label{sec:der}

Let $a>0$, $\omega\neq 0$, $c\in\R$, and
\[
\mathbf r(t)=\bigl(a\cos(\omega t),\,a\sin(\omega t),\,ct\bigr),\qquad t\in\R .
\]

\paragraph{First derivative.}
\[
\mathbf r'(t)=\bigl(-a\omega\sin(\omega t),\,a\omega\cos(\omega t),\,c\bigr).
\]
Hence
\[
|\mathbf r'(t)|^{2}=a^{2}\omega^{2}(\sin^{2}+\cos^{2})+c^{2}
=a^{2}\omega^{2}+c^{2}=:v^{2},
\]
with $v=\sqrt{a^{2}\omega^{2}+c^{2}}>0$. The curve is regular everywhere.

\paragraph{Second and third derivatives.}
\[
\mathbf r''(t)=\bigl(-a\omega^{2}\cos(\omega t),\,-a\omega^{2}\sin(\omega t),\,0\bigr),
\]
\[
\mathbf r'''(t)=\bigl(a\omega^{3}\sin(\omega t),\,-a\omega^{3}\cos(\omega t),\,0\bigr).
\]

\section{Curvature}
\label{sec:k}

The curvature of a regular curve is
\[
\kappa=\frac{|\mathbf r'\times\mathbf r''|}{|\mathbf r'|^{3}}.
\]
The cross product is
\[
\mathbf r'\times\mathbf r''
=\begin{vmatrix}
\mathbf i&\mathbf j&\mathbf k\\
-a\omega\sin & a\omega\cos & c\\
-a\omega^{2}\cos & -a\omega^{2}\sin & 0
\end{vmatrix}
=\bigl(ca\omega^{2}\sin,\,-ca\omega^{2}\cos,\,a^{2}\omega^{3}\bigr).
\]
Its squared norm is
\[
|\mathbf r'\times\mathbf r''|^{2}
=c^{2}a^{2}\omega^{4}+a^{4}\omega^{6}
=a^{2}\omega^{4}(c^{2}+a^{2}\omega^{2})
=a^{2}\omega^{4}v^{2},
\]
so $|\mathbf r'\times\mathbf r''|=a\omega^{2}v$. Therefore
\[
\boxed{\;\kappa=\frac{a\omega^{2}v}{v^{3}}
=\frac{a\omega^{2}}{a^{2}\omega^{2}+c^{2}}\;}
\]
which is constant.

\section{Torsion}
\label{sec:t}

The torsion is
\[
\tau=\frac{(\mathbf r'\times\mathbf r'')\cdot\mathbf r'''}{|\mathbf r'\times\mathbf r''|^{2}}.
\]
The scalar triple product is
\[
(\mathbf r'\times\mathbf r'')\cdot\mathbf r'''
=(ca\omega^{2}\sin)(a\omega^{3}\sin)
+(-ca\omega^{2}\cos)(-a\omega^{3}\cos)
+(a^{2}\omega^{3})(0)
=ca^{2}\omega^{5}(\sin^{2}+\cos^{2})
=ca^{2}\omega^{5}.
\]
Using $|\mathbf r'\times\mathbf r''|^{2}=a^{2}\omega^{4}v^{2}$,
\[
\boxed{\;\tau=\frac{ca^{2}\omega^{5}}{a^{2}\omega^{4}v^{2}}
=\frac{c\omega}{a^{2}\omega^{2}+c^{2}}\;}
\]
which is also constant.

\begin{rem}
The standard helix parametrisation $\mathbf r(t)=(R\cos t,R\sin t,Pt)$ has
$R=a$ and $P=c/\omega$. Then $\kappa=R/(R^{2}+P^{2})$ and
$\tau=P/(R^{2}+P^{2})$, which reduce to the two boxed formulas above.
\end{rem}

\section{Frenet--Serret as a constant-coefficient system}
\label{sec:fs}

Reparametrise by arc length $s=vt$. The Frenet frame
$(\mathbf T,\mathbf N,\mathbf B)$ satisfies
\[
\frac{d}{ds}
\begin{pmatrix}\mathbf T\\ \mathbf N\\ \mathbf B\end{pmatrix}
=
A
\begin{pmatrix}\mathbf T\\ \mathbf N\\ \mathbf B\end{pmatrix},
\qquad
A=
\begin{pmatrix}
0 & \kappa & 0\\
-\kappa & 0 & \tau\\
0 & -\tau & 0
\end{pmatrix},
\]
with constant $\kappa,\tau$. The matrix $A\in\mathfrak{so}(3)$; its
characteristic polynomial is
$\det(A-\lambda I)=-\lambda(\lambda^{2}+\kappa^{2}+\tau^{2})$, so the
eigenvalues are $0,\pm i\Om$ with $\Om=\sqrt{\kappa^{2}+\tau^{2}}$.

\section{Explicit recovery of the curve}
\label{sec:rec}

Because $A$ is constant, $\mathbf T(s)=e^{sA}\mathbf T(0)$. Up to a Euclidean
motion (choice of initial frame and origin) we may align the helix axis with
$\mathbf e_{3}$ and take
\[
\mathbf T(s)=\Bigl(-\frac{\kappa}{\Om}\sin(\Om s),\;
\frac{\kappa}{\Om}\cos(\Om s),\;
\frac{\tau}{\Om}\Bigr).
\]
One checks $|\mathbf T|=1$ and
\[
\mathbf T'(s)=\bigl(-\kappa\cos(\Om s),\,-\kappa\sin(\Om s),\,0\bigr)
=\kappa\,\mathbf N(s),
\]
with $\mathbf N(s)=(-\cos(\Om s),-\sin(\Om s),0)$. The binormal
$\mathbf B=\mathbf T\times\mathbf N$ is
\[
\mathbf B(s)=\Bigl(\frac{\tau}{\Om}\sin(\Om s),\;
-\frac{\tau}{\Om}\cos(\Om s),\;
\frac{\kappa}{\Om}\Bigr),
\]
and a direct computation verifies $\mathbf B'=-\tau\mathbf N$ and
$\mathbf N'=-\kappa\mathbf T+\tau\mathbf B$; thus the Frenet--Serret
equations hold with the prescribed constants.

Integrating $\mathbf r'(s)=\mathbf T(s)$ gives
\[
\boxed{\;
\mathbf r(s)=\frac{\kappa}{\kappa^{2}+\tau^{2}}
\Bigl(\cos(\Om s),\,\sin(\Om s),\,0\Bigr)
+\frac{\tau}{\Om}\,s\,\mathbf e_{3}+\mathbf c\;}
\tag{$\star$}
\]
with constant $\mathbf c$. This is exactly a circular helix of radius
$R=\kappa/(\kappa^{2}+\tau^{2})$ and axis $\mathbf e_{3}$.

\begin{rem}[Common sign/index slip]
A frequently seen but \emph{incorrect} recovery formula writes the axial term
as $(\tau/\kappa)\,s$. That is wrong unless $P=0$: the axial component of
$\mathbf T$ is $\tau/\Om$, not $\tau/\kappa$. Substituting $(\star)$ back into
the curvature and torsion formulas reproduces $\kappa$ and $\tau$
exactly, confirming the corrected bookkeeping.
\end{rem}

\section{Uniqueness and degeneracies}
\label{sec:uniq}

\begin{thm}[Characterisation of constant $\kappa,\tau$ curves]
\label{thm:main}
Let $\gamma:I\to\R^{3}$ be a regular $C^{3}$ curve.
\begin{enumerate}
\item[(i)] If $\kappa\equiv\kappa_{0}>0$ and $\tau\equiv\tau_{0}$ are both
constant, then $\gamma$ is, up to a Euclidean motion, the circular helix
\eqref{star}.
\item[(ii)] If $\kappa\equiv0$ then $\gamma$ is a straight line.
\item[(iii)] If $\kappa>0$ is constant and $\tau\equiv0$ then $\gamma$ is a
circle (planar constant-curvature curve).
\end{enumerate}
\end{thm}

\begin{proof}
The Frenet frame satisfies $\mathbf F'=A\mathbf F$ with the constant matrix
$A$ of \S\ref{sec:fs}. Hence $\mathbf F(s)=e^{sA}\mathbf F(0)$; the tangent
$\mathbf T(s)$ is therefore determined up to the initial frame, i.e.\ up to a
proper rigid motion. Integrating $\mathbf T$ yields $\gamma$ up to a
translation, giving the form \eqref{star}. Cases (ii) and (iii) follow from the
Frenet equations: $\kappa=0$ implies $\mathbf T'=0$ (straight line);
$\tau=0$ gives $\mathbf B'=0$, so the curve lies in the plane orthogonal to
the constant $\mathbf B$ and has constant curvature there, hence a circle.
\end{proof}

\begin{rem}[On terminology]
Lancret's theorem (1811) states the \emph{weaker} characterisation:
$\tau/\kappa=\text{constant}$ iff the tangent makes a constant angle with a
fixed direction (a \emph{general} helix). The present theorem requires
$\kappa$ and $\tau$ \emph{each} to be constant and concludes the stronger
statement that the curve is a \emph{circular} helix. A circular helix is a
special case of a general helix, not the converse.
\end{rem}

\section{Conclusion}
\label{sec:con}

We have proved, with fully explicit derivatives and a correctly booked
recovery formula, that constant curvature together with constant torsion
characterises the circular helix (and its degenerate circle/line limits) among
regular space curves, up to Euclidean motions. The result is a theorem of
classical differential geometry; it entails no physical claim.
\end{document}
